\section{Products and adjacency for the first nonreduced boundary}
\label{sec:products-v163}
The local fibre calculation applies to arbitrary larger Smith exponents
and to any number of distinct contacts. This gives a complete boundary
class, rather than a primitive open in a class whose complement has
not been identified.

Let $R_0$ be a regular symmetric pencil whose line avoids corank at
least three. At its distinct corank-two points write the positive
Smith exponents as $(a_i,b_i)$, $a_i\leq b_i$. Assume
\begin{equation}\label{eq:first-boundary-class-v163}
                      a_i\in\{1,2\}\quad\hbox{for all }i.
\end{equation}
There is no bound on $b_i$ except the dimension constraint on the
pencil. Let $r$ count the points with $a_i=2$ and $s$ those with
$a_i=1$. Denote by $\widehat G$ the normalization of the total reduced
Hilbert graph of lifted pencil lines and by $\mathfrak X_{R_0}$ its
scheme fibre. The polarization on complete quadrics is
$\bigotimes_{q=1}^{n-1}H_q$, with Hilbert polynomial $\binom n2l+1$.

\begin{theorem}[Complete fibres and their component intersections]
\label{thm:product-fibres-v163}
After choosing parameters and residual frames at the distinct contact
points, there is an isomorphism of schemes
\begin{equation}\label{eq:product-fibres-v163}
              \mathfrak X_{R_0}\simeq X_2^{\,r}\times(\Pj^1)^s.
\end{equation}
Frame changes act by the corresponding transition isomorphisms; the
construction is intrinsic without an ordering of the supports.
In particular the reduced fibre has exactly $2^r$ irreducible
components. For a subset $S$ of the $r$ length-two contacts, the
component is
\begin{equation}\label{eq:component-product-v163}
 (\Pj^4)^{|S|}\times(\mathbb F_2)^{r-|S|}\times(\Pj^1)^s,
 \qquad \dim=s+2r+2|S|.
\end{equation}
The intersection of two reduced components is obtained by replacing
each factor where their choices differ by the doubled-line
$\Pj^1$, and leaving the common factors unchanged. These intersections
are reduced. This describes all multiple intersections as well.

The fibre is connected and is reduced exactly when $r=0$. For $r>0$
its nilradical has nilpotence order exactly $r+1$: its $(r+1)$st
power is zero, and its $r$th power is nonzero. Locally its equations
are tensor products of \eqref{eq:full-fibre-ring-v163} and smooth
jet charts. In particular the reduced support is not a replacement
for its Hilbert fibre scheme. The Euler characteristic of the
underlying complex space is $7^r2^s$.
\end{theorem}
\begin{proof}
We construct the morphism before using completed rings. At a
length-two support, the universal curve over $X_2$ contains its
fixed main curve. In the charts \eqref{eq:HB-ideal-v163}, the fibre
relations $eA=eB=e^2C=0$ give
\[
             xG=eH,\qquad xH=-eCG,
             \qquad x^2G=x^2H=0.
\]
Consequently the quotient of the fixed main ideal by the curve ideal
is annihilated by $x^2$. Thus the entire modification, including the
nilpotent directions in its parameter scheme, is specified by a
coherent ideal quotient on the first infinitesimal neighbourhood of
the support. At a length-one contact the corresponding quotient is
annihilated by $x$. These are finite algebraic neighbourhoods, not
unbounded formal data.

Use the central Schur splitting and undo its congruence on these
finite quotients. Only finitely many coefficients of the central
gauge are involved. Include the fixed main lift, and extend by its
forced ideal off the supports. This gives algebraic closed ideal
families in the fixed complete-quadric target. Multiplicative
stability of the ideals is preserved by the target coordinate
isomorphisms. Distinct supports are disjoint, so the constructions
combine independently. Flatness follows from the local flat
families and the patching lemma; their Hilbert polynomial is the
one stated above. Indeed each contact of length $a_i$ removes
$a_i$ from the last exterior degree of the main curve and puts
that degree on the tail. The attachment lengths give Euler
characteristic one for the total curve. The smaller exterior
systems have no base point on this rank open.

This constructs a finite-presentation morphism from the right side
of \eqref{eq:product-fibres-v163} to the fibre of the reduced Hilbert
graph. Each point is a specialization of nonincident graph points:
for primitive tails this is the division family, and for conic tails
it is the dense open $e(A^2+CB^2)\ne0$ in
\eqref{eq:HB-base-v163}. A polynomial line through a specified
point with generic direction meets that open; normalization of
its closure gives an algebraic one-parameter specialization.
The nilpotent part of the constructed map also factors through
the reduced total graph, by the completed comparison in
Proposition~\ref{prop:embedded-functors-v163}, not merely because
its closed points do. That proposition identifies its completed
total neighbourhoods with the regular polynomial graph charts
and the smooth contact factors. Hence the total graph is regular
there, the morphism lifts uniquely to $\widehat G$, and the
completed local maps on the fibre are isomorphisms.

Recovery of the embedded ideal on each disjoint finite neighbourhood
makes this morphism a monomorphism. The formal etaleness criterion
therefore makes it an open immersion. Its source is proper, so it
is also a closed immersion. The fibre of $\widehat G\to G$ is
connected: this is a proper birational morphism from an integral
normal variety to the smooth Grassmannian, and Stein factorization
applies. The nonempty open-and-closed image is the entire fibre.
This proves the scheme isomorphism, not just an equality of reduced
supports. It also explains why the arbitrary values of $b_i$ make
no difference: $k_i$ is removed by a generating-row transformation
which is undone in the fixed incidence target, while the contact
algebra and the resulting finite quotient have length $a_i$.

The reduced components and intersections follow from
Theorem~\ref{thm:complete-fibre-v163} by products. Tensor products
over $\C$ preserve the displayed reduced smooth components and
their reduced intersections. The nilradical is the sum of the
pullbacks $N_i$ of the square-zero nilradicals of the $r$ factors.
The quotient by this sum is reduced, so there is no additional
nilradical. Every monomial of degree $r+1$ repeats a factor and
vanishes. The tensor product of nonzero sections of all $N_i$ is
nonzero on a product of their nonzero affine charts, so the $r$th
power is nonzero. The Euler characteristic follows by multiplicativity
and Corollary~\ref{cor:adjacency-v163}.

Finally the construction uses the actual support schemes and the
fixed embedded ideals. Permuting the support labels permutes the
factors; changing local source frames applies their functorial
transition maps. These satisfy the cocycle identity because the
coordinate changes are undone in the fixed target. This gives the
asserted unlabelled interpretation, rather than a preferred global
ordering of contacts.
\end{proof}

\subsection{Resultant strata on the higher-contact open}
For larger $a_i$ the division chart remains useful even though the
whole proper fibre is not classified here. The following records
exactly its boundary stratification. It is also a check that repeated
roots are not being mistaken for multiple reduced components.

\begin{proposition}[Gcd partitions, codimensions and multiplicity]
\label{prop:resultant-strata-v163}
In $W_a=\{(f,g,w)\}$, with $f$ monic of degree $a$ and
$\deg g,\deg w<a$, fix a partition
$\mu=(\mu_1,\ldots,\mu_k)$ of $d\leq a$. The locally closed locus
where the monic gcd of $f,g$ is
$\prod_{j=1}^k(x-\xi_j)^{\mu_j}$ with distinct $\xi_j$ has dimension
\[
                         3a-2d+k
\]
and codimension $2d-k$. Its finite attaching scheme is the disjoint
union of the primary schemes
$\Spec\C[x]/((x-\xi_j)^{\mu_j})$, and its tails are projective
lines over those schemes. A multiple root is one primary thickening,
not several components. At a point with squarefree $f$ and precisely
$d$ common roots, the resultant divisor is locally a product of
$d$ independent parameters and has multiplicity $d$.
\end{proposition}
\begin{proof}
Parametrize the gcd by its distinct roots, giving $k$ parameters,
and write $f=hf_1$, $g=hg_1$. The monic polynomial $f_1$ has
$a-d$ coefficients and $g_1$ has $a-d$ coefficients; impose the
open coprimality condition on this residual pair and the conditions
specifying the exact gcd partition. Recovering the monic gcd is
regular on the constant-gcd-degree stratum by the fixed-rank
Euclidean linear equations. Up to the finite permutations of equal
parts, this gives dimension $k+2(a-d)$; $w$ contributes $a$.
The statement includes $d=a$, where $g=0$ and $f=h$.
The primary decomposition is the Chinese remainder theorem for
distinct roots and the gcd fibre equation of
Theorem~\ref{thm:division-chart-v162}. At a squarefree $f$, its
roots are etale local coordinates and
$\operatorname{Res}(f,g)=\prod_{f(\xi)=0}g(\xi)$.
The evaluations of a polynomial of degree $<a$ at these roots are
independent by the Vandermonde matrix. The nonzero factors are
units, proving the final assertion. No normal-crossing statement
is made at a multiple root of $f$.
\end{proof}

\subsection{A testable extension beyond length two}
The following conjecture concerns the reduced components of a local
fibre, not the miniversal congruence quotient. It distinguishes two
questions: component geometry and a possible nilpotent structure on
their union.

\begin{remark}[Conjecture: rational components of the corank-two contact fibre]
\label{conj:contact-components-v163}
For every $a\geq1$, let $X_a$ be the fibre over $(x^a,0,0)$ of the
normalization of the total reduced graph $\Gamma_a\to B_a$ defined
in Theorem~\ref{thm:division-chart-v162}. Every irreducible component
of $(X_a)_{\rm red}$ has rational normalization. The closure of the
primitive jet open is an irreducible component of dimension $a$,
and its intersections with the other components contain a locus of
curves whose tail cycle has a repeated line through the attaching
point. For $a\geq2$ the scheme $X_a$ is not reduced.
\end{remark}
For $a=1$ the fibre is $\Pj^1$. For $a=2$ the conjecture follows
from the complete description above, including the doubled-line
intersection and the square-zero ideal. For $a\geq3$ it is a
conjecture, not a component classification. The statement is
falsifiable separately by the birational type of a component, its
adjacency to the primitive component, or reducedness of the fibre.
An eventual classification should refine it by the degrees and
primary thicknesses of tail cycles of total degree $a$. The present
product theorem makes no assertion about higher-corank contacts or
singular pencils in the Hilbert problem; the independent inverse
and singular-pencil invariant theorems continue to apply on their
original domains.
